\setcounter{section}{5}
\section{Metrados vinculados al diseño}
Conductor de terna tendida=3\ensuremath{\times}\ensuremath{\Sigma}Lc\_\allowbreak{}EDS=20,274.361 m. Puentes/colas provisionales=3fases\ensuremath{\times}ramales retenidos\ensuremath{\times}2 m=960.000 m. Reserva3\% sobre ambos=638 m. Compra de estudio redondeada=21873 m. El2 m por ramal es una provisión que debe sustituirse por la longitud del detalle de armado, diferenciada de la reserva de suministro.

\begingroup\small\setlength{\LTpre}{4pt}\setlength{\LTpost}{6pt}
\begin{longtable}{*{3}{L{\dimexpr\linewidth/3-2\tabcolsep\relax}}}
\toprule
\textbf{Recurso} & \textbf{Unidad} & \textbf{Cantidad de estudio} \\ \midrule
\endfirsthead
\toprule
\textbf{Recurso} & \textbf{Unidad} & \textbf{Cantidad de estudio} \\ \midrule
\endhead
Armado ANCLAJE & jgo & 22.000 \\
Armado ANCLAJE\_\allowbreak{}PROPUESTO & jgo & 16.000 \\
Armado ANGULO & jgo & 41.000 \\
Armado SUSPENSION & jgo & 30.000 \\
Armado TERMINAL & jgo & 2.000 \\
Conductor AAAC 50: terna tendida & m & 20,274.361 \\
Conductor para puentes/colas & m & 960.000 \\
Poste CAC 13 m / 600 daN tipo & ud & 72.000 \\
Poste CAC 15 m / 600 daN tipo & ud & 30.000 \\
Poste CAC 17 m / 600 daN tipo & ud & 9.000 \\
Reserva adicional conductor & m & 638.000 \\
Retenida y anclaje por ramal & ud & 160.000 \\
Transformador 250 kVA netos en sitio & ud & 2.000 \\
\bottomrule\end{longtable}\endgroup

metrado\_\allowbreak{}detallado.csv conserva ID, recurso, fórmula, poste/vano y estado. Los civiles no reciben volumen cerrado: excavación=\ensuremath{\pi}D\textsuperscript{2}h/4 y concreto=excavación-volumen desplazado cuando se defina geometría real. Un bloque 1.2 m o una malla dibujada no acredita capacidad ni resistividad. Las planillas dejan esas dimensiones vacías.

